% Part: first-order-logic
% Chapter: model-theory
% Section: reducts-and-expansions

\documentclass[../../../include/open-logic-section]{subfiles}

\begin{document}

\olfileid{mod}{bas}{red}
\section{न्यूनीकृत आणि विस्तारित रचना}

ज्या भाषांत काही चिन्हे समान आहेत अशा भाषांची, तसेच त्या भाषांसाठीच्या
!!{structure}s ची, तुलना करणे अनेकदा उपयुक्त किंवा आवश्यक असते. सर्वांत
नेहमीचे प्रकरण असे असते की !!a{language}~$\Lang{L}$ मधील सर्व चिन्हे
!!a{language}~$\Lang{L'}$ मध्येही असतात, म्हणजे $\Lang{L} \subseteq
\Lang{L'}$. मग अतिरिक्त चिन्हांची अर्थनिर्धारणे जोडून आणि समान चिन्हांची
अर्थनिर्धारणे तशीच ठेवून $\Lang{L}$-!!{structure}~$\Struct{M}$ चा विस्तार
नेहमीच $\Lang{L'}$-!!{structure} मध्ये करता येतो. दुसरीकडे,
$\Lang{L}$ मध्ये नसलेल्या चिन्हांची अर्थनिर्धारणे केवळ ``विसरून''
$\Lang{L'}$-रचना~$\Struct{M'}$ पासून $\Lang{L}$-रचना मिळवता येते.

\begin{defn}
\ollabel{defn:reduct}
समजा $\Lang L \subseteq \Lang L'$, $\Struct M$ ही
$\Lang L$-!!{structure} आणि $\Struct M'$ ही $\Lang L'$-!!{structure} आहे.
$\Struct M$ ही $\Struct M'$ ची $\Lang L$ पर्यंतची \emph{न्यूनीकृत रचना}
आणि $\Struct M'$ ही $\Struct M$ ची $\Lang L'$ पर्यंतची
\emph{विस्तारित रचना} असेल तेव्हा आणि केवळ तेव्हाच पुढील अटी पूर्ण होतात:
\begin{enumerate}
\item $\Domain{M} = \Domain{M'}$
\item प्रत्येक !!{constant}~$c \in \Lang L$ साठी $\Assign{c}{M} =
  \Assign{c}{M'}$.
\item प्रत्येक !!{function}~$f \in \Lang L$ साठी $\Assign{f}{M} =
  \Assign{f}{M'}$.
\item प्रत्येक !!{predicate}~$P \in \Lang L$ साठी $\Assign{P}{M} =
  \Assign{P}{M'}$.
\end{enumerate}
\end{defn}

\begin{prop}
\ollabel{prop:reduct}
एखादी $\Lang{L}$-!!{structure}~$\Struct{M}$ ही एखाद्या
$\Lang{L'}$-!!{structure}~$\Struct{M'}$ ची न्यूनीकृत रचना असेल, तर सर्व
$\Lang{L}$-!!{sentence}s~$!A$ साठी,
\[
\Sat{M}{!A} \text{ iff } \Sat{M'}{!A}.
\]
\end{prop}

\begin{proof}
  सराव.
\end{proof}

\begin{prob}
\olref[mod][bas][red]{prop:reduct} सिद्ध करा.
\end{prob}

\begin{defn}
आपल्याकडे $\Lang{L}$-रचना $\Struct{M}$ असेल, एका $n$-स्थानी
!!{predicate}~$P$ ची भर घालून मिळालेला $\Lang{L}$ चा विस्तार
$\Lang{L'} = \Lang{L} \cup \{P\}$ असेल, आणि $R \subseteq \Domain{M}^n$
हा $n$-स्थानी संबंध असेल, तर $\Assign{P}{M'} = R$ असलेली
$\Struct{M}$ ची विस्तारित रचना~$\Struct{M'}$ दर्शवण्यासाठी आपण
$\Expan{M}{R}$ असे लिहितो.
\end{defn}

\end{document}
